\chapter{Cứng và Mềm}
\markboth{Cứng và Mềm}{Firmness and Softness}

\section{Cứng trong nguyên tắc, mềm trong cách đối xử.}
\begin{truyen}{Cứng trong nguyên tắc, mềm trong cách đối xử.}{Principles can be firm while manners stay gentle.}
\chuhoa{M}{ột thầy giáo không chấp nhận gian lận, nhưng khi học sinh vi phạm, thầy không sỉ nhục trước lớp. Thầy phân tích rõ, xử lý nghiêm, rồi vẫn chừa cho em một lối quay lại. Cái cứng đáng quý nhất là cái cứng không làm mất nhân phẩm người khác.}
\textit{(A teacher did not tolerate cheating, yet when a student violated the rule he did not humiliate the child publicly. He explained clearly, handled the matter firmly, and still left a path back. The most valuable firmness is the one that does not strip others of dignity.)}
\end{truyen}
\begin{baihoc}
Cứng về điều đúng không đồng nghĩa với thô về cách làm.
\textit{(Being firm about what is right does not require being rough in action.)}
\end{baihoc}
\begin{ghinhoanh}
\textbf{Short English to remember:}

Firm truth can wear a gentle face.
\end{ghinhoanh}
\begin{tuhocnhanh}
1. Em thường thiên về cứng quá hay mềm quá? \textit{(Do you tend to be too hard or too soft?)}

2. Em có thể giữ nguyên tắc mà vẫn tử tế hơn bằng cách nào? \textit{(How can you keep your principles while being kinder?)}
\end{tuhocnhanh}
\ngancach

\section{Mềm quá nên ai cũng dễ lấn.}
\begin{truyen}{Mềm quá nên ai cũng dễ lấn.}{Softness without boundaries invites harm.}
\chuhoa{C}{ô sợ làm mất lòng nên việc gì cũng nhận, yêu cầu nào cũng cố chiều. Lâu dần người khác quen với việc đẩy phần khó sang cho cô. Sự hiền lành không đi cùng ranh giới rõ ràng rất dễ biến thành cửa mở cho sự lợi dụng.}
\textit{(She feared upsetting others, so she accepted every task and tried to please every request. Over time people grew used to pushing difficult parts onto her. Kindness without clear boundaries easily becomes an open gate for exploitation.)}
\end{truyen}
\begin{baihoc}
Mềm không có nghĩa là không biết nói “không”.
\textit{(Being gentle does not mean being unable to say no.)}
\end{baihoc}
\begin{ghinhoanh}
\textbf{Short English to remember:}

Kindness needs boundaries.
\end{ghinhoanh}
\begin{tuhocnhanh}
1. Em có đang chiều người khác đến mức mệt mình không? \textit{(Are you pleasing others to the point of exhausting yourself?)}

2. Em cần đặt ranh giới rõ hơn ở đâu? \textit{(Where do you need clearer boundaries?)}
\end{tuhocnhanh}
\ngancach

\section{Cứng đầu khác với kiên định.}
\begin{truyen}{Cứng đầu khác với kiên định.}{Stubbornness is not the same as conviction.}
\chuhoa{A}{nh tưởng mình bản lĩnh vì không đổi ý. Nhưng thật ra nhiều lần anh chỉ không chịu nghe thêm vì sợ mất mặt. Kiên định là biết giữ điều đúng sau khi đã lắng nghe; còn cứng đầu là bịt tai để giữ cái tôi.}
\textit{(He thought he was strong because he never changed his mind. Yet many times he simply refused to listen further because he feared losing face. Conviction means holding what is right after listening; stubbornness means closing the ears to protect the ego.)}
\end{truyen}
\begin{baihoc}
Cái cứng có trí tuệ luôn đi cùng khả năng học hỏi.
\textit{(Intelligent firmness always goes together with the ability to learn.)}
\end{baihoc}
\begin{ghinhoanh}
\textbf{Short English to remember:}

Conviction listens; stubbornness refuses.
\end{ghinhoanh}
\begin{tuhocnhanh}
1. Em đang kiên định hay đang cứng đầu ở vấn đề nào? \textit{(In which issue are you being convicted, and in which are you merely stubborn?)}

2. Ý kiến nào em nên nghe thêm trước khi khẳng định mình đúng? \textit{(Which perspective should you hear more fully before insisting you are right?)}
\end{tuhocnhanh}
\ngancach

\section{Mềm lời để cứu một cuộc nói chuyện khó.}
\begin{truyen}{Mềm lời để cứu một cuộc nói chuyện khó.}{A soft tone can keep truth alive.}
\chuhoa{K}{hi con thú nhận điểm kém, người cha định quát lớn. Nhưng rồi ông hạ giọng, hỏi kỹ, và nghe ra đằng sau điểm số là cả một quãng mất phương hướng. Nhờ mềm lại ở giọng nói, cuộc nói chuyện không biến thành bức tường.}
\textit{(When the child admitted a poor score, the father first wanted to shout. Yet he lowered his voice, asked carefully, and heard behind the score an entire season of confusion. Because the tone softened, the conversation did not turn into a wall.)}
\end{truyen}
\begin{baihoc}
Lời mềm không làm sự thật yếu đi; nó giúp sự thật có chỗ đi vào.
\textit{(A soft tone does not weaken the truth; it helps the truth enter.)}
\end{baihoc}
\begin{ghinhoanh}
\textbf{Short English to remember:}

Soft words can carry hard truth.
\end{ghinhoanh}
\begin{tuhocnhanh}
1. Có cuộc trò chuyện nào của em sẽ tốt hơn nếu giọng nói mềm lại? \textit{(Which conversation of yours would improve if your tone softened?)}

2. Em có đang nhầm giữa nghiêm túc và nặng nề không? \textit{(Are you confusing seriousness with harshness?)}
\end{tuhocnhanh}
\ngancach

\section{Biết lúc nào cần cứng để bảo vệ điều yếu hơn.}
\begin{truyen}{Biết lúc nào cần cứng để bảo vệ điều yếu hơn.}{Firmness can protect the vulnerable.}
\chuhoa{K}{hi thấy bạn nhỏ hơn bị bắt nạt, cô chị không chọn im cho êm chuyện. Cô đứng ra chặn lại, nói rõ giới hạn, và báo với người có trách nhiệm. Có lúc mềm là tốt, nhưng cũng có lúc phải cứng để điều yếu hơn không bị nghiền nát.}
\textit{(When she saw a younger student being bullied, the older sister did not stay quiet for the sake of peace. She stepped in, set boundaries clearly, and informed a responsible adult. There are moments when gentleness is good, but there are also moments when firmness is needed so that the vulnerable are not crushed.)}
\end{truyen}
\begin{baihoc}
Cái cứng đẹp nhất là cái cứng đứng về phía người yếu thế.
\textit{(The most beautiful firmness is the kind that stands with the vulnerable.)}
\end{baihoc}
\begin{ghinhoanh}
\textbf{Short English to remember:}

Firmness can defend the weak.
\end{ghinhoanh}
\begin{tuhocnhanh}
1. Em có điều gì cần mạnh lên để bảo vệ đúng điều đúng? \textit{(What do you need to become stronger in so that you can protect what is right?)}

2. Khi nào sự mềm của em vô tình tiếp tay cho điều xấu? \textit{(When might your softness unintentionally enable something wrong?)}
\end{tuhocnhanh}
\ngancach
